
\subsection{Drift detection}
\label{sec:bg:drift}

Data drift.



\subsection{Hyperplane-Assisted Softmax}
\label{sec:bg:has}


\subsubsection{Motivation and Key Insight}

HASeparator \cite{kansizoglou-santavas-etal:2020} was introduced as an
alternative to centre-based angular loss functions such as ArcFace
\cite{deng-jiankang-etal:2019}. Unlike methods that implicitly treat the last-layer
weight vector of each class as that class's prototype centre, HASeparator
explicitly operates on the separation hyperplanes between pairs
of classes.
The separation hyperplane between class $i$ and class $j$ is uniquely
defined by the characteristic vector
$\bar{\bw}_{ij} = \bar{\bw}_i - \bar{\bw}_j$
\cite{kansizoglou-bampis-gasteratos:2022}, making its computation exact and
free of the approximation error inherent in centre-based methods.
The method adapts the maximum-margin principle of Support Vector
Machines~(SVM) \cite{cortes-vapnik:1995} to the continuously evolving
embedding space of a CNN, where both feature vectors and decision
boundaries are jointly optimised.

\subsubsection{Forward Pass and Logit Computation}

Let $\bE \in \Rn^{B \times N}$ denote a batch of feature embeddings,
$\bW \in \Rn^{N \times C}$ the classification layer weights,
$B$ the batch size, $N = 64$ the embedding dimension, and
$C$ the number of classes.
Both embeddings and weights are $\ell_2$-normalised, yielding
$\hat{\bE}$ and $\hat{\bW}$, so that their matrix product corresponds
to cosine similarities.
The scaled logits are
\begin{equation}
  O^i_j \;=\; \sigma\,\hat{E}^i_k\,\hat{W}^k_j \,,
  \label{eq:logits}
\end{equation}
where $\sigma > 0$ is a learnable feature scaler that defines the
radius of the embedding hypersphere, and index notation follows the
Einstein summation convention \cite{kansizoglou-santavas-etal:2020}.

\subsubsection{HASeparator Penalty}

For each sample $i$ with ground-truth label $\ell_i$, the
$C$ separation hyperplanes of the target class are formed by
broadcasting the target weight vector $\hat{\bw}_{\ell_i}$
across all class weight vectors:
\begin{equation}
  \bH \;=\; \hat{\bW}^{+} - \bT^{+} \;\in\; \Rn^{B \times N \times C} \,,
  \label{eq:hyperplanes}
\end{equation}
where $\hat{\bW}^{+}$ is $\hat{\bW}$ replicated $B$ times and
$\bT^{+}$ holds $\hat{\bw}_{\ell_i}$ replicated across the $C$
class axis.
After $\ell_2$-normalizing $\bH$ along the neuron axis to obtain
$\hat{\bH}$, the projection of each embedding onto the $C$
hyperplane normals is
\begin{equation}
  P^i_j \;=\; \hat{E}^{\cdot i}_k\,\hat{H}^{\cdot i k}_j \,,
  \label{eq:projections}
\end{equation}
producing cosine values bounded in $[-1,1]$.
A relaxed SVM cost is then applied element-wise:
\begin{equation}
  J(p^i_j) \;=\; m - \min(p^i_j,\, m)\,, \quad m \in (0,1]\,,
  \label{eq:cost}
\end{equation}
and the total batch penalty is
\begin{equation}
  C_t \;=\; \frac{1}{B}\,\bm{1}_i\,J^i_j\,\bm{1}^j \,.
  \label{eq:ct}
\end{equation}
Summing over all $C$ hyperplanes of the target class penalises embeddings
that lie close to \emph{any} decision boundary, directly maximising the
minimum angular distance to the nearest hyperplane.

\subsubsection{Training Objective and Warmup Schedule}

The final objective combines $C_t$ with the negative log-likelihood loss:
\begin{equation}
  \mathcal{L} \;=\; \alpha\,C_{\mathrm{NLL}} \;+\; \beta_{\mathrm{eff}}(t)\,C_t \,,
  \label{eq:total_loss}
\end{equation}
where $\alpha$ and $\beta$ are weighting coefficients.
To prevent the angular penalty from destabilising the embedding geometry
in early training, $\beta$ is linearly warmed up from zero:
\begin{equation}
  \beta_{\mathrm{eff}}(t) \;=\; \beta \cdot \min\!\left(\frac{t}{T_w},\,1\right)\,,
  \label{eq:warmup}
\end{equation}
with $T_w = 5$ warm-up epochs.
Convergence health is monitored through the mean pairwise cosine
similarity of the $\ell_2$-normalised class weight vectors,
$\bar{w}_{\cos}$; values in $[0.10,\,0.25]$ indicate well-separated
prototypes, while values approaching unity signal geometric collapse.



